\documentclass[11pt,a4paper]{article}
\usepackage[utf8]{inputenc}
\usepackage[T1]{fontenc}
\usepackage[brazil]{babel}
\usepackage{lmodern}
\usepackage[margin=2cm]{geometry}
\usepackage{amsmath,amssymb}
\usepackage{booktabs}
\usepackage{graphicx}
\usepackage{hyperref}
\usepackage{xcolor}
\usepackage{titlesec}
\usepackage{tabularx}
\usepackage{enumitem}
\usepackage{float}
\usepackage{array}
\renewcommand{\tabularxcolumn}[1]{>{\raggedright\arraybackslash}p{#1}}
\usepackage{newunicodechar}
\usepackage[most]{tcolorbox}
\usepackage{etoolbox}

\hypersetup{colorlinks=true, linkcolor=blue!70!black, citecolor=blue!70!black, urlcolor=blue!70!black}

\titleformat{\section}{\large\bfseries}{\thesection.}{0.5em}{}
\titleformat{\subsection}{\normalsize\bfseries}{\thesubsection.}{0.5em}{}
\titlespacing*{\section}{0pt}{1.2ex plus 0.5ex minus .2ex}{0.6ex}
\titlespacing*{\subsection}{0pt}{1.0ex plus 0.4ex minus .2ex}{0.4ex}
\setlength{\parskip}{0.45em}
\setlength{\parindent}{0pt}
\setlist{itemsep=0.15em,topsep=0.3em}

% --- símbolos ---
\newunicodechar{π}{\ensuremath{\pi}}
\newunicodechar{τ}{\ensuremath{\tau}}
\newunicodechar{φ}{\ensuremath{\varphi}}
\newunicodechar{ψ}{\ensuremath{\psi}}
\newunicodechar{Σ}{\ensuremath{\Sigma}}
\newunicodechar{→}{\ensuremath{\rightarrow}}
\newunicodechar{⇒}{\ensuremath{\Rightarrow}}
\newunicodechar{∈}{\ensuremath{\in}}
\newunicodechar{−}{\ensuremath{-}}
\newunicodechar{≈}{\ensuremath{\approx}}
\newunicodechar{≠}{\ensuremath{\neq}}
\newunicodechar{≡}{\ensuremath{\equiv}}
\newunicodechar{≤}{\ensuremath{\leq}}
\newunicodechar{≥}{\ensuremath{\geq}}
\newunicodechar{′}{\ensuremath{'}}
\newunicodechar{‑}{-}
\newunicodechar{✔}{\ensuremath{\checkmark}}
\newunicodechar{✘}{\ensuremath{\times}}

% --- cores e caixas ---
\definecolor{cC}{HTML}{7A4FA3}
\definecolor{cS}{HTML}{2F7D74}
\definecolor{cAcc}{HTML}{2F5D8A}
\newcommand{\metrictag}[2]{{\setlength{\fboxsep}{1.5pt}%
  \ifstrequal{#1}{c}{\colorbox{cC}{\textcolor{white}{\scriptsize\textbf{#2}}}}{%
  \ifstrequal{#1}{s}{\colorbox{cS}{\textcolor{white}{\scriptsize\textbf{#2}}}}{%
  \colorbox{black!60}{\textcolor{white}{\scriptsize\textbf{#2}}}}}}}
\newcommand{\pagechip}[1]{{\setlength{\fboxsep}{1.5pt}\colorbox{cAcc}{\textcolor{white}{\scriptsize\texttt{#1}}}}}
\newcommand{\readtag}[1]{\texttt{\small #1}}
\newcommand{\partbar}[1]{\par\vspace{1.2em}\noindent\colorbox{black!80}{\parbox{\dimexpr\linewidth-2\fboxsep}{\color{white}\bfseries\sffamily\MakeUppercase{#1}}}\par\vspace{0.4em}}
\newtcolorbox{citacard}[1]{breakable, enhanced, colback=white, colframe=black!25, boxrule=0.5pt, arc=2pt,
  left=6pt, right=6pt, top=4pt, bottom=4pt, title={\small #1}, coltitle=black, colbacktitle=black!6, fonttitle=\normalfont}
\newenvironment{origquote}{\par\itshape\small}{\par}
\newenvironment{tradquote}{\par\small\textsc{Tradução:}\ }{\par}
\newcommand{\explab}[1]{\par\smallskip{\small\color{cAcc}\textbf{\MakeUppercase{#1}}}\par}
\colorlet{notebg}{blue!4}\colorlet{notefg}{cAcc}
\colorlet{okbg}{green!5}\colorlet{okfg}{green!45!black}
\colorlet{warnbg}{orange!8}\colorlet{warnfg}{orange!80!black}
\newtcolorbox{infobox}[2][note]{breakable, enhanced, boxrule=0pt, arc=2pt, left=6pt, right=6pt, top=4pt, bottom=4pt,
  colback=#1bg, borderline west={3pt}{0pt}{#1fg}, frame hidden,
  before upper={\ifstrempty{#2}{}{\textbf{#2}\par}}}
\colorlet{pairc}{cC}\colorlet{pairs}{cS}
\newenvironment{paircol}[2]{\begin{tcolorbox}[width=0.485\linewidth, nobeforeafter, box align=top, boxrule=0.5pt,
  colframe=black!25, colback=white, borderline north={3pt}{0pt}{pair#1}, left=5pt, right=5pt,
  before upper={{\color{pair#1}\small\textbf{\MakeUppercase{#2}}}\par}]}{\end{tcolorbox}}

\newunicodechar{ℓ}{\ensuremath{\ell}}
\newunicodechar{⊙}{\ensuremath{\odot}}
\newunicodechar{α}{\ensuremath{\alpha}}
\newunicodechar{β}{\ensuremath{\beta}}

\begin{document}

\begin{center}
    {\LARGE \textbf{Fichamento: limiar que maximiza o F1}} \\[0.35em]
    {\large LIPTON, Z. C.; ELKAN, C.; NARYANASWAMY, B. \textit{Optimal Thresholding of Classifiers to Maximize F1 Measure}. In: Machine Learning and Knowledge Discovery in Databases (ECML PKDD), 2014, p.~225--239} \\[0.3em]
    \textbf{Tema:} relação entre o melhor F1 e o limiar de decisão, em problemas binários e multirrótulo \\[0.2em]
    \textbf{Aluno:} Arthur Santos de Oliveira (RA: 156297) \\[0.2em]
    \textbf{Atualizado em:} 07/11/2026 \\[0.2em]
    \textbf{Link:} \url{https://arxiv.org/abs/1402.1892} \\[0.2em]
    \textbf{Leitura:} seções 1 a 4 da versão arXiv (título \textit{Thresholding Classifiers to Maximize F1 Score}, 16 páginas); estudo de caso lido pelo resumo \\[0.2em]
    \textbf{Repositório:} \url{https://github.com/Arthur-so/shap-ecg-arrhythmia}
\end{center}
\vspace{-0.5em}
\hrule
\vspace{0.8em}

\textbf{Como ler.} \texttt{[p. N]} indica a página do PDF do arXiv. As citações estão no original em inglês, cada uma seguida de uma tradução livre e de uma explicação.

\begin{infobox}[note]{Resumo}
\begin{itemize}
\item \textbf{Objetivo:} entender como escolher o limiar de decisão para maximizar o F1, em classificação binária e multirrótulo (F1 micro, macro e por exemplo).
\item \textbf{Resultado teórico:} para qualquer classificador com saída real, o limiar ótimo está ligado ao melhor F1 alcançável; se as probabilidades são bem calibradas, o limiar ótimo é a \textbf{metade} do F1 ótimo.
\item \textbf{Alertas:} com um classificador sem informação, o F1 é maximizado prevendo tudo como positivo; o F1-macro dá, na prática, mais peso às classes raras.
\end{itemize}
\end{infobox}

\section{O que o trabalho faz}
Os autores derivam a relação entre o limiar de decisão e o F1 que ele produz, e mostram que o limiar que maximiza o F1 não é, em geral, 0,5 [p.~1]. No caso calibrado, o limiar ótimo é $F1^*/2$ [p.~1]. Discutem também como os três modos de média do F1 no caso multirrótulo afetam o limiar de cada rótulo, em especial dos rótulos raros [p.~1--2].

\section{Citações}

\begin{citacard}{\pagechip{p. 1} Resumo}
\begin{origquote}“For any classifier that produces a real-valued output, we derive the relationship between the best achievable F1 score and the decision-making threshold that achieves this optimum. As a special case, if the classifier outputs are well-calibrated conditional probabilities, then the optimal threshold is half the optimal F1 score.”\end{origquote}
\begin{tradquote}“Para qualquer classificador com saída real, derivamos a relação entre o melhor F1 alcançável e o limiar de decisão que o alcança. Como caso particular, se as saídas do classificador são probabilidades condicionais bem calibradas, o limiar ótimo é a metade do F1 ótimo.”\end{tradquote}
\explab{Explicação}
Justifica escolher o limiar de cada classe pelo F1, em vez de fixar 0,5. Também ajuda a ler os limiares do TCC: se o modelo fosse calibrado, os limiares ficariam perto de $F1/2$ (0,4 a 0,5); valores muito acima disso, como 0,99 no IAVB, indicam que o modelo dá probabilidades altas a ECGs sem a classe, isto é, não é calibrado nessa classe.
\end{citacard}

\begin{citacard}{\pagechip{p. 1} Classificador sem informação}
\begin{origquote}“As another special case, if the classifier is completely uninformative, then the optimal behavior is to classify all examples as positive. Since the actual prevalence of positive examples typically is low, this behavior can be considered undesirable.”\end{origquote}
\begin{tradquote}“Como outro caso particular, se o classificador não tem nenhuma informação, o comportamento ótimo é classificar todos os exemplos como positivos. Como a prevalência real dos positivos costuma ser baixa, esse comportamento pode ser indesejável.”\end{tradquote}
\explab{Explicação}
É um cuidado do procedimento: maximizar o F1 na validação pode levar a limiares muito baixos em classes que o modelo não aprendeu. No TCC isso não ocorreu (o menor limiar foi 0,34, em STD), mas o alerta justifica conferir os limiares.
\end{citacard}

\begin{citacard}{\pagechip{p. 2} F1-macro e classes raras}
\begin{origquote}“In this sense, macro average F1 can be argued not to weigh labels equally, but actually to give greater weight to performance on rare labels.”\end{origquote}
\begin{tradquote}“Nesse sentido, pode-se argumentar que o F1-macro não pesa os rótulos igualmente, mas dá, na verdade, mais peso ao desempenho nos rótulos raros.”\end{tradquote}
\explab{Explicação}
O F1-macro é o critério da parada antecipada no TCC; classes raras como STE e BRE pesam nele tanto quanto BRD, e as suas variações entre épocas movem mais o F1-macro.
\end{citacard}

\section{Relação com o TCC}
Dá a base teórica do limiar por classe escolhido na validação (seção 1.2.2 do texto), que no TCC segue o código de Zhang et al.\ (2021). O artigo não trata de ECG nem de explicações.

\section{Como citar no TCC}
\begin{itemize}
\item ``O limiar que maximiza o F1 depende da classe e, em geral, não é 0,5 (Lipton, Elkan e Naryanaswamy, 2014); por isso, ele foi escolhido para cada classe na validação.''
\end{itemize}

\end{document}
